\par\noindent\textbf{(1)} 质心速度 V=mv/(m+M)，相对初速 u=v。入射粒子在质心系速率 v\_* = Mv/(m+M)。散射后其质心系速度为 v\_*(cosχ,sinχ)，故实验室速度 v'\_x=V+v\_*cosχ、v'\_y=v\_*sinχ，靶的反冲由总动量守恒确定。
\par\noindent\textbf{(2)} 实验室角满足 tanθ\_lab=v'\_y/v'\_x=sinχ/(cosχ+m/M)，须用 atan2 保留正确象限。粒子实验室速度模 |v'|=v√[m²+M²+2mMcosχ]/(m+M)，与靶反冲动能一起满足弹性总能量守恒。
\par\noindent\textbf{(3)} 若 m=M，则 tanθ=tan(χ/2)，在 0≤χ≤π 的正向分支有 θ=χ/2，故入射粒子不后向反弹。若 m>M，视实验室终速度方向可见 v'\_x>0；取 dθ/dχ=0 得 cosχ=−M/m，从而 sinθ\_max=M/m，即 θ\_max=arcsin(M/m)<π/2。
\par\noindent\textbf{(4)} 由于散射轴对称，角度映射给 (dσ/dΩ)\_lab=Σ\_branches(dσ/dΩ)\_cm |sinχ/(sinθ\_lab)·dχ/dθ\_lab|。当 m>M 的 θ\_lab 小于最大值时可能有两个质心角映至同一实验室角，必须对分支求和；在转折点 Jacobian 发散属几何caustic，不意味着总散射概率无穷。
